% mazzi: 01
\chapter{Le fonti del diritto internazionale in materia di doping}

% ---------------------------------------------------------------------------
\section{Il quadro normativo internazionale}

Il \textbf{legislatore internazionale} in materia di antidoping si articola su due linee, che
producono fonti distinte.

\begin{center}
\begin{forest} classalbero
[Legislatore internazionale
  [Consiglio d'Europa \\ UNESCO
    [Convenzione Strasburgo 1989]
    [Convenzione Parigi 2005]
  ]
  [WADA
    [Codice mondiale e standard]
  ]
]
\end{forest}
\end{center}

In \textit{ambito internazionale} la materia dell'antidoping è disciplinata da:
\begin{enumerate}
  \item \textbf{Convenzione contro il doping di Strasburgo}: adottata il 16 novembre 1989 e
    ratificata in Italia con la legge 29 novembre 1995, n.\,522;
  \item \textbf{Convenzione internazionale contro il doping UNESCO}: adottata a Parigi dalla
    XXXIII Conferenza generale UNESCO il 19 ottobre 2005 e ratificata in Italia con la legge 26
    novembre 2007, n.\,230;
  \item \textbf{Codice WADA} (\textit{World Anti-Doping Agency}): emanato nel 2004, aggiornato
    nel 2009 e nel 2015;
  \item \textbf{Standard internazionali WADA}.
\end{enumerate}

% ---------------------------------------------------------------------------
\section{La Convenzione di Strasburgo}

\textbf{Convenzione contro il doping di Strasburgo}: adottata il 16 novembre 1989 e ratificata
in Italia con la legge 29 novembre 1995, n.\,522. Persegue due obiettivi:
\begin{itemize}
  \item \textbf{coordinamento interno}, tra istituzioni pubbliche e sportive;
  \item \textbf{cooperazione internazionale}.
\end{itemize}

% ---------------------------------------------------------------------------
\section{La Convenzione UNESCO}

\textbf{Convenzione internazionale contro il doping UNESCO}: adottata a Parigi dalla XXXIII
Conferenza generale UNESCO il 19 ottobre 2005 e ratificata in Italia con la legge 26 novembre
2007, n.\,230; sottoscritta da 191 paesi partecipanti alla Conferenza Generale di Parigi.

\rubrica{Contenuti principali}
\begin{itemize}
  \item misure più incisive rispetto alla Convenzione di Strasburgo;
  \item riconoscimento della \textbf{leadership della WADA};
  \item lista WADA e procedure di \textbf{Esenzioni a Fini Terapeutici} allegate;
  \item \textbf{Codice WADA} come appendice alla Convenzione stessa;
  \item adozione mondiale del Codice WADA.
\end{itemize}

\rubrica{Ripartizione delle responsabilità}
\begin{itemize}
  \item la \textbf{comunità sportiva} è impegnata nei regolamenti, nei controlli antidoping,
    nonché nelle sanzioni e nelle misure disciplinari;
  \item i \textbf{governi} sono i maggiori responsabili delle misure legislative (in materia di
    tutela sanitaria), delle misure finanziarie e dei laboratori (compreso il campo della
    ricerca); coordinano inoltre le politiche e l'operato dei servizi e degli organismi pubblici
    impegnati nella lotta antidoping, e incoraggiano le organizzazioni sportive ad elaborare e
    applicare le misure appropriate di loro competenza;
  \item comunità sportiva e governi si ripartiscono la responsabilità delle \textbf{misure
    educative} e dei \textbf{programmi di informazione}.
\end{itemize}

% ---------------------------------------------------------------------------
\section{Il Codice WADA}

\textbf{Codice WADA}: insieme delle norme applicate a livello internazionale che regolamenta la
materia del doping e dell'antidoping.
\begin{itemize}
  \item armonizza le politiche, le norme e i regolamenti antidoping dei singoli paesi e
    chiarisce le responsabilità dei firmatari;
  \item obbliga all'accettazione della \textbf{Lista delle sostanze e metodi proibiti}.
\end{itemize}

\rubrica{Punti salienti}
\begin{itemize}
  \item definizione di doping \textbf{più restrittiva e più precisa};
  \item lista di sostanze;
  \item puntualizzazione per i controlli:
    \begin{itemize}
      \item ordinari, senza preavviso, in e fuori competizione;
      \item elenco mirato di atleti da sottomettere a controlli mirati (\textbf{TDP});
    \end{itemize}
  \item utilizzo di \textbf{standard internazionali} uguali:
    \begin{itemize}
      \item per tutti i laboratori;
      \item per le sostanze soggette a restrizione;
      \item per le modalità di controllo.
    \end{itemize}
\end{itemize}

\rubrica{Soggetti tenuti al rispetto del Codice}
Le regole e i principi dell'antidoping devono essere seguiti da:
\begin{itemize}
  \item \textbf{Comitati Olimpici Nazionali};
  \item \textbf{Federazioni Sportive Internazionali};
  \item \textbf{Organismi di controllo nazionali} (\textbf{NADO}).
\end{itemize}

% ---------------------------------------------------------------------------
\section{Gli standard internazionali}

Il Codice WADA rimanda a sei standard internazionali, ciascuno dedicato a un ambito specifico:
\begin{itemize}
  \item \textbf{Lista delle sostanze e metodi proibiti} (\textit{The Prohibited List});
  \item \textbf{Controlli ed investigazioni} (\textit{International Standard for Testing and
    Investigations});
  \item \textbf{Esenzioni ai fini terapeutici} (\textit{International Standard for TUE});
  \item \textbf{Laboratori} (\textit{International Standard for Laboratory});
  \item \textbf{Privacy} (\textit{International Standard for Protection of Privacy and Personal
    Information});
  \item \textbf{Compliance} (\textit{International Standard For Code Compliance By
    Signatories}).
\end{itemize}
